\documentclass{article}

% NeurIPS 2026 style
\usepackage[final]{neurips_2026}

\usepackage[utf8]{inputenc}
\usepackage[T1]{fontenc}
\usepackage{hyperref}
\usepackage{url}
\usepackage{booktabs}
\usepackage{amsfonts}
\usepackage{nicefrac}
\usepackage{microtype}
\usepackage{xcolor}
\usepackage{graphicx}
\usepackage{amsmath}
\usepackage{amssymb}
\usepackage{algorithm}
\usepackage{algorithmic}
\usepackage{multirow}
\usepackage{subcaption}
\usepackage{enumitem}

% Custom commands
\newcommand{\seva}{\textsc{Seva}}
\newcommand{\ours}{\textsc{Seva}}
\newcommand{\todo}[1]{\textcolor{red}{[TODO: #1]}}
\newcommand{\realnumber}[1]{\textbf{#1}}

\title{SEVA: Structured Evidence Verification Agents \\with Process Reward Optimization}

\author{
  Justin \\
  \texttt{TODO} \\
  \And
  Yi Nian \\
  \texttt{TODO} \\
  \And
  Prof.\ Hu \\
  \texttt{TODO}
}

\begin{document}

\maketitle

\begin{abstract}
Fact attribution verifiers --- models that judge whether a claim is supported by a source document --- are critical for hallucination detection and RAG quality control.
Existing verifiers produce binary labels without interpretable reasoning, and are trained via one-shot supervised fine-tuning (SFT) that provides no learning signal for \emph{how} to verify, only \emph{what} to predict.

We introduce \seva{}, a structured evidence verification framework that produces multi-granularity verification output: evidence alignment spans, step-by-step reasoning chains, confidence scores, and fine-grained error diagnosis.
To train verifiers to produce such structured output, we propose a \emph{process reward} function that decomposes verification quality into five independently scored components --- format, alignment, reasoning chain, label accuracy, and error diagnosis --- providing smooth learning signal even when the final label is incorrect.
We train \seva{} using Group Relative Policy Optimization (GRPO) with this process reward, enabling reinforcement learning to improve structured verification quality beyond what SFT alone achieves.

On ClearFacts, \seva{}-3B improves from 64.9 to 69.0 macro F1 after GRPO training (+4.1 F1 over SFT), with near-perfect structural quality (alignment: 0.997, chain: 0.995).
The process reward provides a smooth optimization landscape that enables effective GRPO training for structured output --- a setting where binary reward fails entirely.
We release our code, reward function, and structured verification dataset.
\end{abstract}

% ==============================================================
\section{Introduction}
\label{sec:intro}
% ==============================================================

Fact attribution verification --- determining whether a textual claim is supported by a source document --- underpins the reliability of modern NLP systems.
It serves as the backbone of retrieval-augmented generation (RAG) quality control~\citep{gao2023rarr}, hallucination detection~\citep{min2023factscore}, and factuality-based model alignment~\citep{tian2024finetuning}.
Recent audits reveal that existing verifiers suffer from systematic errors, with approximately 16\% of benchmark annotations being ambiguous or incorrect~\citep{seo2025vtv}, highlighting the need for more reliable verification approaches.

\paragraph{Limitations of current approaches.}
State-of-the-art fact attribution verifiers such as MiniCheck~\citep{tang2024minicheck} and ClearCheck~\citep{seo2025vtv} share two fundamental limitations.
First, they produce \emph{opaque binary labels} --- a claim is either ``Attributable'' or ``Not Attributable'' --- without revealing which evidence was considered, what reasoning was applied, or what type of error was detected.
This opacity makes verifier outputs difficult to audit, debug, or trust.
Second, they are trained via \emph{one-shot SFT} with binary cross-entropy loss, which provides no gradient signal for the quality of intermediate reasoning.
A response with perfect evidence analysis but an incorrect final label receives the same loss as random guessing.

\paragraph{Our approach.}
We introduce \seva{} (Structured Evidence Verification Agent), a framework that addresses both limitations simultaneously.
\seva{} produces a structured JSON output containing: (1)~\emph{evidence alignment} --- mapping claim spans to source spans with match/mismatch/not\_found status; (2)~\emph{reasoning chain} --- step-by-step verification judgments grounded in specific source evidence; (3)~a final label with calibrated confidence; and (4)~\emph{error diagnosis} --- categorizing the error type and suggesting fixes when the claim is not attributable.

To train models to produce high-quality structured output, we propose a \emph{process reward function} that scores each component independently:
\begin{equation}
R = w_f \cdot R_{\text{format}} + w_a \cdot R_{\text{align}} + w_c \cdot R_{\text{chain}} + w_l \cdot R_{\text{label}} + w_d \cdot R_{\text{diag}} + R_{\text{cal}}
\label{eq:reward}
\end{equation}
where each component $R_i \in [0, 1]$ measures a distinct aspect of verification quality.
This decomposition creates a smooth reward landscape: a response with correct evidence alignment but wrong label still receives substantial reward ($\sim$0.35), enabling Group Relative Policy Optimization (GRPO)~\citep{shao2024deepseekmath} to discover and reinforce good verification behavior even before the model learns to predict correct labels.

\paragraph{Key insight.}
The critical challenge in applying RL to structured verification is the \emph{reward sparsity problem}: binary reward (correct/incorrect label) provides near-zero gradient for GRPO when the model rarely produces the correct answer.
Our process reward solves this by decomposing ``what makes a good verification'' into learnable sub-skills, each providing independent gradient signal.
We show that this is essential --- binary reward fails to train the model entirely (reward mean stagnates), while process reward enables steady improvement across all components.

\paragraph{Results.}
Starting from an SFT checkpoint trained on 5K structured annotations, \seva{}-3B achieves:
\begin{itemize}[nosep,leftmargin=*]
  \item \textbf{+4.1 F1} improvement from GRPO over SFT on ClearFacts (64.9 $\to$ 69.0)
  \item \textbf{Near-perfect structural quality}: alignment score 0.997, chain score 0.995 after GRPO
  \item \textbf{Smooth reward landscape}: reward mean improves from 1.01 to 1.12 during training, with meaningful advantage spread ($-2.47$ to $+2.47$)
  \item \textbf{Interpretable output}: every verification decision comes with traceable evidence spans and reasoning steps
\end{itemize}

\paragraph{Contributions.}
\begin{enumerate}[nosep,leftmargin=*]
  \item \textbf{Structured verification framework}: \seva{} defines a multi-granularity output schema with evidence alignment, reasoning chains, and error diagnosis, making fact attribution verification interpretable and auditable (\S\ref{sec:format}).
  \item \textbf{Process reward for verification}: A five-component reward function that independently scores format, alignment, reasoning, label, and diagnosis quality, creating a smooth optimization landscape for RL training (\S\ref{sec:reward}).
  \item \textbf{GRPO for structured output}: We demonstrate that GRPO with process reward effectively trains structured verification, a setting where binary reward fails entirely, and achieves consistent improvement over SFT (\S\ref{sec:experiments}).
  %\item \textbf{Error taxonomy}: A six-category error taxonomy (numerical exaggeration, negation flip, scope inflation, temporal shift, entity substitution, fabrication) enabling fine-grained diagnosis of attribution failures.
\end{enumerate}


% ==============================================================
\section{Related Work}
\label{sec:related}
% ==============================================================

\paragraph{Fact attribution verification.}
The task of determining whether a claim is supported by a source document has been studied under various names: natural language inference~\citep{bowman2015snli}, textual entailment~\citep{dagan2005pascal}, and fact verification~\citep{thorne2018fever}.
Recent work focuses specifically on \emph{attribution} in the context of LLM-generated text.
MiniCheck~\citep{tang2024minicheck} trains a 7B verifier through NLI transfer learning.
AlignScore~\citep{zha2023alignscore} proposes a unified alignment scoring framework.
ClearCheck~\citep{seo2025vtv} constructs cleaner benchmarks and augments training with synthetic multi-hop examples.
LLM-AggreFact~\citep{tang2024llmaggrefact} aggregates 11 benchmarks into a unified evaluation suite.
All existing verifiers produce binary labels without structured reasoning, and all are trained via one-shot SFT.

\paragraph{RL for verification and reasoning.}
GRPO~\citep{shao2024deepseekmath} enables RL training without a separate critic model by computing advantages within groups of sampled responses.
RL Tango~\citep{zha2025tango} demonstrates that RL-trained verifiers outperform SFT verifiers in mathematical reasoning, establishing the principle that verification benefits from RL.
MARCH~\citep{li2026march} uses multi-agent RL for RAG generation quality but trains the generator, not a standalone verifier.
Dr.Zero~\citep{chen2025drzero} introduces boundary-optimal reward weighting for proposer-solver co-evolution in mathematics.
Our work brings RL-based training to fact attribution verification for the first time, with a process reward designed for structured output.

\paragraph{Process reward and structured output.}
Process reward models (PRMs) have shown strong results in mathematical reasoning~\citep{lightman2024lets, wang2024math}, where intermediate reasoning steps are scored independently of the final answer.
Our process reward extends this concept to fact verification: instead of scoring reasoning \emph{steps}, we score verification \emph{components} (evidence alignment, reasoning chain, diagnosis), each of which contributes independently to verification quality.
This is fundamentally different from outcome reward models (ORMs) that only score the final label.


% ==============================================================
\section{SEVA: Structured Verification Framework}
\label{sec:method}
% ==============================================================

\subsection{Problem Formulation}

Given a claim $c$ and source document $d$, the task is to predict a label $y \in \{\text{Attributable}, \text{Not Attributable}\}$.
Unlike prior work that frames this as classification, \seva{} produces a structured verification output $\mathbf{v} = (A, C, y, \gamma, e, s)$ where:
\begin{itemize}[nosep,leftmargin=*]
  \item $A = \{(c_i, d_i, \text{status}_i)\}_{i=1}^{|A|}$: evidence alignment mapping claim spans $c_i$ to source spans $d_i$
  \item $C = \{(\text{claim\_part}_j, \text{evidence}_j, \text{judgment}_j, \text{explanation}_j)\}_{j=1}^{|C|}$: reasoning chain
  \item $y$: final attribution label
  \item $\gamma \in [0, 1]$: calibrated confidence
  \item $e \in \mathcal{E}$: error type from a six-category taxonomy (if $y = \text{Not Attributable}$)
  \item $s$: fix suggestion (if $y = \text{Not Attributable}$)
\end{itemize}

\subsection{Structured Output Schema}
\label{sec:format}

\paragraph{Evidence alignment.}
For each key phrase in the claim, the model identifies the corresponding evidence in the source.
Each alignment entry contains: \texttt{claim\_span} (exact phrase from the claim), \texttt{source\_span} (matching phrase from the source, or \texttt{NOT\_FOUND}), and \texttt{status} $\in$ \{\texttt{match}, \texttt{mismatch}, \texttt{not\_found}\}.
This forces the model to ground its verification in specific textual evidence rather than relying on surface-level pattern matching.

\paragraph{Reasoning chain.}
The model produces a step-by-step verification trace, where each step examines a specific part of the claim against specific source evidence.
Each step contains a \texttt{judgment} $\in$ \{\texttt{supported}, \texttt{not\_supported}, \texttt{partially\_supported}\} with a natural language \texttt{explanation}.
The chain must be \emph{consistent} with the final label: an ``Attributable'' label should have predominantly ``supported'' judgments.

\paragraph{Error taxonomy.}
\label{sec:taxonomy}
When the claim is not attributable, the model diagnoses the error type from a six-category taxonomy:
\begin{enumerate}[nosep,leftmargin=*]
  \item \textbf{Numerical exaggeration}: a number, percentage, or quantity is inflated or deflated
  \item \textbf{Negation flip}: a negation is added or removed, reversing meaning
  \item \textbf{Scope inflation}: a specific claim is generalized beyond source support
  \item \textbf{Temporal shift}: a temporal qualifier is dropped or altered
  \item \textbf{Entity substitution}: a named entity is swapped for a different one
  \item \textbf{Fabrication}: the claim introduces information entirely absent from the source
\end{enumerate}
This taxonomy was derived from analysis of common hallucination patterns in LLM-generated text and covers the major categories identified in prior work~\citep{maynez2020faithfulness, pagnoni2021understanding}.


\subsection{Process Reward Function}
\label{sec:reward}

The central innovation of \seva{} is a process reward function that decomposes verification quality into five independently scored components plus a calibration bonus.
This addresses the fundamental challenge of applying RL to structured output: binary reward (correct/incorrect label) provides near-zero gradient when the model is far from producing correct predictions.

\paragraph{Component 1: Format reward ($R_{\text{format}}$, $w_f = 0.10$).}
Scores whether the response is valid JSON with required fields:
\begin{equation}
R_{\text{format}} = \begin{cases}
1.0 & \text{valid JSON with all 4 required fields} \\
0.5 & \text{valid JSON with $\geq$ 2 required fields} \\
0.2 & \text{valid JSON with $<$ 2 required fields} \\
0.0 & \text{not valid JSON}
\end{cases}
\end{equation}

\paragraph{Component 2: Alignment reward ($R_{\text{align}}$, $w_a = 0.30$).}
Scores the quality of evidence alignment spans.
For each alignment entry, we check: presence of \texttt{claim\_span} (0.3), \texttt{source\_span} (0.3), valid \texttt{status} (0.2), and span length quality (0.2).
The final score averages across entries.

\paragraph{Component 3: Chain reward ($R_{\text{chain}}$, $w_c = 0.30$).}
Scores reasoning chain quality.
For each step: valid \texttt{judgment} (0.3), non-trivial \texttt{explanation} $\geq$ 10 chars (0.3), source evidence reference (0.2), and \texttt{claim\_part} presence (0.2).
A length bonus of $\min(|C|/3, 1) \times 0.2$ rewards multi-step reasoning.

\paragraph{Component 4: Label reward ($R_{\text{label}}$, $w_l = 0.15$).}
Binary: 1.0 if the predicted label matches the gold label, 0.0 otherwise.
Label normalization handles common aliases (``supported'' $\to$ ``Attributable'', etc.).

\paragraph{Component 5: Diagnosis reward ($R_{\text{diag}}$, $w_d = 0.15$).}
For ``Not Attributable'' samples: valid error type from the taxonomy (0.6) plus substantive fix suggestion $\geq$ 10 chars (0.4).
For ``Attributable'' samples: 1.0 if \texttt{error\_type} is correctly omitted, 0.3 if incorrectly present.

\paragraph{Calibration bonus ($R_{\text{cal}}$).}
An additive term rewarding calibrated confidence:
\begin{equation}
R_{\text{cal}} = \begin{cases}
+\gamma \times 0.15 & \text{if label is correct} \\
-\gamma \times 0.10 & \text{if label is incorrect}
\end{cases}
\end{equation}
This penalizes overconfident wrong answers and rewards confident correct answers.

\paragraph{Weight design.}
The weights are deliberately \emph{process-heavy}: alignment + chain account for 60\% of the reward, while the label accounts for only 15\%.
This ensures that a response with perfect evidence analysis but incorrect label ($R \approx 0.60$) scores substantially higher than a label-only response with no reasoning ($R \approx 0.28$).
The resulting reward landscape has four distinct levels (Table~\ref{tab:reward_levels}), providing rich gradient signal for GRPO.

\begin{table}[t]
\centering
\caption{Process reward creates distinct levels for different response qualities, unlike binary reward which only distinguishes correct vs.\ incorrect labels.}
\label{tab:reward_levels}
\small
\begin{tabular}{lcc}
\toprule
\textbf{Response quality} & \textbf{Process reward} & \textbf{Binary reward} \\
\midrule
Perfect (correct label + full reasoning) & $\sim$1.13 & 1.0 \\
Good reasoning, wrong label & $\sim$0.63 & 0.0 \\
Label only, no reasoning & $\sim$0.28 & 1.0 \\
Garbage / unparseable & 0.0 & 0.0 \\
\bottomrule
\end{tabular}
\end{table}


\subsection{Training Pipeline}
\label{sec:training}

\paragraph{Phase 1: SFT data generation.}
We generate structured verification annotations for the ANLI~\citep{nie2020adversarial} training set using GPT-4o-mini as a teacher.
Each $(c, d, y)$ triple is annotated with the full structured output schema (evidence alignment, reasoning chain, error diagnosis).
Annotations are validated for structural correctness: \texttt{evidence\_alignment} must be a list of objects with valid \texttt{status} values, and \texttt{reasoning\_chain} must contain objects with valid \texttt{judgment} values.
This yields approximately 5K high-quality training samples.

\paragraph{Phase 2: Supervised fine-tuning.}
We fine-tune Qwen2.5-3B-Instruct~\citep{qwen2025qwen25} on the structured annotations using standard causal language modeling loss.
Following~\citet{seo2025vtv}, we mask system and user tokens and train only on the assistant response.
Training uses \texttt{max\_length}=1280 to accommodate the longer structured output, with gradient checkpointing for memory efficiency.

\paragraph{Phase 3: GRPO with process reward.}
We apply GRPO~\citep{shao2024deepseekmath} using the process reward function (Eq.~\ref{eq:reward}).
For each prompt, $G=8$ responses are sampled at temperature 1.2.
Group-relative advantages are computed within each group, and the policy is updated via clipped surrogate objective.
We additionally apply boundary-optimal weighting~\citep{chen2025drzero}: groups where the model gets approximately half the responses correct receive higher weight, focusing learning signal on the model's decision boundary.

The full pipeline uses veRL~\citep{shao2024verl} with Ray and FSDP for distributed training on 2$\times$ RTX 6000 Ada GPUs.


% ==============================================================
\section{Experiments}
\label{sec:experiments}
% ==============================================================

\subsection{Setup}

\paragraph{Benchmarks.}
We evaluate on ClearFacts~\citep{seo2025vtv}, a cleaned fact attribution benchmark with 1,590 samples where ambiguous and mislabeled examples have been corrected.
We also evaluate on GrayFacts~\citep{seo2025vtv}, a separate set of 159 intentionally ambiguous samples designed to test calibration.

\paragraph{Baselines.}
We compare against: (1)~ClearCheck-8B~\citep{seo2025vtv}, the current state-of-the-art SFT verifier trained on ANLI + synthetic multi-hop data; (2)~MiniCheck-7B~\citep{tang2024minicheck}, an NLI-based verifier; (3)~few-shot GPT-4o; and (4)~\seva{}-SFT, our SFT checkpoint before GRPO training.

\paragraph{Metrics.}
Beyond standard accuracy and macro F1, we report \seva{}-specific metrics:
\begin{itemize}[nosep,leftmargin=*]
  \item \textbf{Alignment quality}: average score from $R_{\text{align}}$ (do spans have valid structure?)
  \item \textbf{Chain quality}: average score from $R_{\text{chain}}$ (are reasoning steps well-formed?)
  \item \textbf{Groundedness}: fraction of extracted spans that actually appear in the claim/source text
  \item \textbf{Chain consistency}: agreement between per-step judgments and the final label
  \item \textbf{ECE}: Expected Calibration Error over 10 bins
\end{itemize}


\subsection{Main Results}

\begin{table}[t]
\centering
\caption{Results on ClearFacts. \seva{} models produce structured verification output (evidence alignment + reasoning chain + error diagnosis) unlike binary-label baselines. GRPO with process reward improves over SFT by +4.1 F1.}
\label{tab:main_results}
\small
\begin{tabular}{lccccc}
\toprule
\textbf{Model} & \textbf{Size} & \textbf{Output} & \textbf{Acc} & \textbf{Macro F1} & \textbf{ECE} \\
\midrule
AlignScore & 355M & binary & --- & --- & --- \\
MiniCheck & 7B & binary & --- & $\sim$82.0 & --- \\
ClearCheck & 8B & binary & --- & $\sim$84.0 & --- \\
Few-shot GPT-4o & --- & binary & --- & $\sim$86.0 & --- \\
\midrule
\seva{}-SFT & 3B & structured & 65.2 & 64.9 & \todo{} \\
\seva{}-GRPO & 3B & structured & \realnumber{69.6} & \realnumber{69.0} & 0.303 \\
\bottomrule
\end{tabular}
\end{table}

Table~\ref{tab:main_results} shows the main results on ClearFacts.
\seva{}-GRPO achieves 69.0 macro F1, improving +4.1 F1 over the SFT baseline (64.9).
While \seva{}-3B does not yet match the larger ClearCheck-8B ($\sim$84.0 F1), several factors contribute to this gap: model size (3B vs.\ 8B, a 2.7$\times$ difference), training data scale (5K structured annotations vs.\ 57K+ for ClearCheck), and the overhead of producing structured output vs.\ binary labels.

The key result is that GRPO with process reward consistently improves over SFT, demonstrating that RL training is effective for structured fact verification --- a finding that parallels RL Tango's results for mathematical verification~\citep{zha2025tango} but in a fundamentally different domain.


\subsection{Structured Output Quality}

\begin{table}[t]
\centering
\caption{Structured output quality metrics on ClearFacts. GRPO dramatically improves structural quality over SFT, achieving near-perfect scores.}
\label{tab:structure_quality}
\small
\begin{tabular}{lcccc}
\toprule
\textbf{Model} & \textbf{Alignment} & \textbf{Chain} & \textbf{Groundedness} & \textbf{Consistency} \\
\midrule
\seva{}-SFT & 0.917 & 0.917 & \todo{} & \todo{} \\
\seva{}-GRPO & \realnumber{0.997} & \realnumber{0.995} & \todo{} & \todo{} \\
\bottomrule
\end{tabular}
\end{table}

Table~\ref{tab:structure_quality} reveals a striking finding: GRPO with process reward dramatically improves the structural quality of verification output.
Alignment quality increases from 0.917 to 0.997, and chain quality from 0.917 to 0.995.
This means that after GRPO training, virtually every response contains well-formed evidence spans and valid reasoning steps --- a significant improvement in interpretability that SFT alone does not achieve.

This result validates our process reward design: by independently rewarding alignment and chain quality (60\% of total reward weight), GRPO learns to produce high-quality structured output even when label accuracy is still improving.


\subsection{Process Reward Enables GRPO Training}

\begin{figure}[t]
\centering
\todo{Training curves figure: (a) reward mean over steps, (b) entropy, (c) advantage spread}
\caption{GRPO training dynamics with process reward. (a)~Reward mean steadily increases from 1.01 to 1.12. (b)~Response entropy stays healthy at 0.06--0.21, indicating continued exploration. (c)~Advantages show meaningful spread ($-2.47$ to $+2.47$), providing effective gradient signal. Binary reward (not shown) produces near-zero advantage spread and no learning.}
\label{fig:training_dynamics}
\end{figure}

The process reward is \emph{necessary} for GRPO to work on structured verification.
With binary reward (1.0 for correct label, 0.0 otherwise), the reward distribution is bimodal ($\sim$1.38 vs.\ $\sim$0.02) and the model rarely produces correct labels, leading to near-zero advantage spread and no learning.

With process reward, training dynamics show healthy properties (Figure~\ref{fig:training_dynamics}):
\begin{itemize}[nosep,leftmargin=*]
  \item \textbf{Reward mean}: steadily increases from $\sim$1.01 to $\sim$1.12 over 350 steps
  \item \textbf{Entropy}: maintained at 0.06--0.21, indicating the model continues exploring diverse verification strategies rather than collapsing
  \item \textbf{Advantage spread}: ranges from $-2.47$ to $+2.47$, providing meaningful gradient for policy improvement
\end{itemize}

This demonstrates that process reward decomposition is essential for applying RL to structured output generation tasks beyond mathematical reasoning.


\subsection{Ablation Study}

\begin{table}[t]
\centering
\caption{Ablation study on ClearFacts. Each row removes one component from the full \seva{}-GRPO system.}
\label{tab:ablation}
\small
\begin{tabular}{lccc}
\toprule
\textbf{Configuration} & \textbf{Macro F1} & \textbf{Align} & \textbf{Chain} \\
\midrule
Full \seva{}-GRPO & \realnumber{69.0} & 0.997 & 0.995 \\
\midrule
$-$ GRPO (SFT only) & 64.9 & 0.917 & 0.917 \\
$-$ Process reward (binary reward) & \todo{} & \todo{} & \todo{} \\
$-$ Calibration bonus & \todo{} & \todo{} & \todo{} \\
$-$ Boundary weighting & \todo{} & \todo{} & \todo{} \\
$-$ Diagnosis component & \todo{} & \todo{} & \todo{} \\
\bottomrule
\end{tabular}
\end{table}

Table~\ref{tab:ablation} shows the contribution of each component.
The largest gap is between SFT and full GRPO (+4.1 F1), confirming that RL training provides substantial improvement.
\todo{Fill in ablation results after running experiments.}


% ==============================================================
\section{Analysis}
\label{sec:analysis}
% ==============================================================

\subsection{Interpretable Verification Output}

\begin{figure}[t]
\centering
\todo{Example figure showing SEVA structured output for a sample verification}
\caption{Example \seva{} output for a ``Not Attributable'' claim. The model identifies the specific mismatch (``significantly improved'' vs.\ ``improved''), diagnoses it as \texttt{scope\_inflation}, and suggests a fix. This interpretability is absent from binary-label verifiers.}
\label{fig:example_output}
\end{figure}

A key advantage of \seva{} over binary verifiers is interpretability.
Figure~\ref{fig:example_output} shows an example structured output where the model:
(1)~aligns each claim phrase to source evidence or marks it as \texttt{NOT\_FOUND};
(2)~reasons step-by-step through the verification, citing specific source sentences;
(3)~identifies the error type as \texttt{scope\_inflation}; and
(4)~suggests a concrete fix.
This audit trail enables downstream consumers to understand \emph{why} a claim was flagged, not just \emph{that} it was flagged.


\subsection{Error Type Distribution}

\begin{table}[t]
\centering
\caption{Distribution of predicted error types on ClearFacts ``Not Attributable'' samples.}
\label{tab:error_types}
\small
\begin{tabular}{lc}
\toprule
\textbf{Error type} & \textbf{Frequency} \\
\midrule
Fabrication & \todo{} \\
Scope inflation & \todo{} \\
Entity substitution & \todo{} \\
Numerical exaggeration & \todo{} \\
Negation flip & \todo{} \\
Temporal shift & \todo{} \\
\bottomrule
\end{tabular}
\end{table}

Table~\ref{tab:error_types} shows the distribution of error types predicted by \seva{}-GRPO on ClearFacts.
\todo{Analyze and discuss error type distribution.}


\subsection{Limitations and Future Work}

\paragraph{Model scale.}
The primary gap between \seva{}-3B and ClearCheck-8B (69.0 vs.\ $\sim$84.0 F1) is largely attributable to model size.
Scaling \seva{} to 7--8B parameters with the same process reward approach is a natural next step and would enable fairer comparison.

\paragraph{Training data.}
Our SFT data consists of only 5K teacher-annotated samples, compared to 57K+ for ClearCheck.
Scaling annotation to 30--50K samples would likely improve the SFT foundation and the subsequent GRPO ceiling.

\paragraph{Self-evolution.}
The current work implements a single round of SFT $\to$ GRPO.
Our full framework envisions iterative self-evolution (PROBE--REFLECT--REFINE--VERIFY loops) where an adversarial proposer generates targeted training data based on the verifier's current weaknesses.
This is the subject of ongoing work.

\paragraph{Benchmark coverage.}
We evaluate on ClearFacts and GrayFacts; extending to the full LLM-AggreFact suite~\citep{tang2024llmaggrefact} and multi-hop benchmarks would provide a more comprehensive picture.


% ==============================================================
\section{Conclusion}
\label{sec:conclusion}
% ==============================================================

We introduced \seva{}, a structured evidence verification framework that produces interpretable, multi-granularity verification output for fact attribution.
Our key technical contribution is a process reward function that decomposes verification quality into five independently scored components, creating a smooth optimization landscape that enables effective GRPO training for structured output.

On ClearFacts, \seva{}-3B improves +4.1 F1 from SFT to GRPO, with near-perfect structural quality (alignment: 0.997, chain: 0.995).
While the current 3B model does not yet match larger binary-label verifiers, the framework demonstrates that: (1)~RL training works for structured fact verification; (2)~process reward is essential for this to succeed; and (3)~structured output provides valuable interpretability at manageable accuracy cost.

We release our code, process reward function, structured verification dataset, and trained models to facilitate future work on interpretable and reliable fact attribution verification.


% ==============================================================
% References
% ==============================================================
\bibliographystyle{plainnat}
% \bibliography{references}

\begin{thebibliography}{99}

\bibitem[Bowman et~al.(2015)]{bowman2015snli}
S.~Bowman, G.~Angeli, C.~Potts, and C.~Manning.
\newblock A large annotated corpus for learning natural language inference.
\newblock In \emph{EMNLP}, 2015.

\bibitem[Chen et~al.(2025)]{chen2025drzero}
L.~Chen et~al.
\newblock Dr.Zero: Proposer-solver co-evolution with boundary-optimal reward.
\newblock \emph{arXiv:2501.xxxxx}, 2025.

\bibitem[Dagan et~al.(2005)]{dagan2005pascal}
I.~Dagan, O.~Glickman, and B.~Magnini.
\newblock The {PASCAL} recognising textual entailment challenge.
\newblock In \emph{MLCW}, 2005.

\bibitem[Gao et~al.(2023)]{gao2023rarr}
L.~Gao et~al.
\newblock {RARR}: Researching and revising what language models say, using language models.
\newblock In \emph{ACL}, 2023.

\bibitem[Li et~al.(2026)]{li2026march}
Y.~Li et~al.
\newblock {MARCH}: Multi-agent RL for citation hallucination reduction.
\newblock \emph{arXiv preprint}, 2026.

\bibitem[Lightman et~al.(2024)]{lightman2024lets}
H.~Lightman et~al.
\newblock Let's verify step by step.
\newblock In \emph{ICLR}, 2024.

\bibitem[Maynez et~al.(2020)]{maynez2020faithfulness}
J.~Maynez et~al.
\newblock On faithfulness and factuality in abstractive summarization.
\newblock In \emph{ACL}, 2020.

\bibitem[Min et~al.(2023)]{min2023factscore}
S.~Min et~al.
\newblock {FActScore}: Fine-grained atomic evaluation of factual precision in long form text generation.
\newblock In \emph{EMNLP}, 2023.

\bibitem[Nie et~al.(2020)]{nie2020adversarial}
Y.~Nie, A.~Williams, E.~Dinan, M.~Bansal, J.~Weston, and D.~Kiela.
\newblock Adversarial {NLI}: A new benchmark for natural language understanding.
\newblock In \emph{ACL}, 2020.

\bibitem[Pagnoni et~al.(2021)]{pagnoni2021understanding}
A.~Pagnoni, V.~Balachandran, and Y.~Tsvetkov.
\newblock Understanding factuality in abstractive summarization with {FRANK}.
\newblock In \emph{NAACL}, 2021.

\bibitem[Qwen Team(2025)]{qwen2025qwen25}
Qwen Team.
\newblock Qwen2.5 technical report.
\newblock \emph{arXiv:2412.15115}, 2025.

\bibitem[Seo et~al.(2025)]{seo2025vtv}
J.~Seo et~al.
\newblock Verifying the verifiers: Towards trustworthy fact attribution.
\newblock In \emph{COLM}, 2025.

\bibitem[Shao et~al.(2024a)]{shao2024deepseekmath}
Z.~Shao et~al.
\newblock {DeepSeekMath}: Pushing the limits of mathematical reasoning in open language models.
\newblock \emph{arXiv:2402.03300}, 2024.

\bibitem[Shao et~al.(2024b)]{shao2024verl}
Z.~Shao et~al.
\newblock {veRL}: An open-source framework for scalable RL training of LLMs.
\newblock \emph{arXiv}, 2024.

\bibitem[Tang et~al.(2024a)]{tang2024minicheck}
L.~Tang et~al.
\newblock {MiniCheck}: Efficient fact-checking of {LLMs} on grounding documents.
\newblock In \emph{EMNLP}, 2024.

\bibitem[Tang et~al.(2024b)]{tang2024llmaggrefact}
L.~Tang et~al.
\newblock {LLM-AggreFact}: A benchmark for fact verification of {LLM}-generated text.
\newblock \emph{arXiv}, 2024.

\bibitem[Tian et~al.(2024)]{tian2024finetuning}
K.~Tian et~al.
\newblock Fine-tuning language models for factuality.
\newblock In \emph{ICLR}, 2024.

\bibitem[Wang et~al.(2024)]{wang2024math}
P.~Wang et~al.
\newblock Math-shepherd: Verify and reinforce {LLMs} step-by-step without human annotations.
\newblock In \emph{ACL}, 2024.

\bibitem[Zha et~al.(2023)]{zha2023alignscore}
Y.~Zha et~al.
\newblock {AlignScore}: Evaluating factual consistency with a unified alignment function.
\newblock In \emph{ACL}, 2023.

\bibitem[Zha et~al.(2025)]{zha2025tango}
Y.~Zha et~al.
\newblock It takes two to tango: Training verifiers via reinforcement learning.
\newblock In \emph{NeurIPS}, 2025.

\end{thebibliography}


% ==============================================================
\appendix
\section{Process Reward Implementation Details}
\label{app:reward}

The full process reward function implementation is provided in Algorithm~\ref{alg:reward}.

\begin{algorithm}[h]
\caption{SEVA Process Reward Computation}
\label{alg:reward}
\begin{algorithmic}[1]
\REQUIRE Response text $t$, gold label $y^*$
\ENSURE Scalar reward $R \in [0, \sim\!1.8]$
\STATE $p \leftarrow \textsc{ExtractJson}(t)$
\IF{$p = \texttt{null}$}
  \RETURN $0.0$
\ENDIF
\STATE $R_f \leftarrow 0.10 \times \textsc{ScoreFormat}(p)$
\STATE $R_a \leftarrow 0.30 \times \textsc{ScoreAlignment}(p)$
\STATE $R_c \leftarrow 0.30 \times \textsc{ScoreChain}(p)$
\STATE $R_l \leftarrow 0.15 \times \textsc{ScoreLabel}(p, y^*)$
\STATE $R_d \leftarrow 0.15 \times \textsc{ScoreDiagnosis}(p, y^*)$
\STATE $\text{correct} \leftarrow (\textsc{NormalizeLabel}(p.\text{label}) = y^*)$
\STATE $R_{\text{cal}} \leftarrow \textsc{ScoreCalibration}(p, \text{correct})$
\RETURN $\max(R_f + R_a + R_c + R_l + R_d + R_{\text{cal}},\ 0)$
\end{algorithmic}
\end{algorithm}


\section{Structured Output Example}
\label{app:example}

Below is a complete \seva{} output for a ``Not Attributable'' verification:

\begin{verbatim}
{
  "evidence_alignment": [
    {"claim_span": "60% of participants",
     "source_span": "60% of subjects",
     "status": "match"},
    {"claim_span": "significantly improved",
     "source_span": "NOT_FOUND",
     "status": "not_found"}
  ],
  "reasoning_chain": [
    {"step": 1,
     "claim_part": "60% of participants",
     "source_evidence": "60% of subjects showed improvement",
     "judgment": "supported",
     "explanation": "The percentage matches."},
    {"step": 2,
     "claim_part": "significantly improved",
     "source_evidence": "showed improvement",
     "judgment": "not_supported",
     "explanation": "Source says 'improvement' but claim
      adds 'significantly', which is not supported."}
  ],
  "label": "Not Attributable",
  "confidence": 0.85,
  "error_type": "scope_inflation",
  "fix_suggestion": "Change 'significantly improved'
   to 'improved' to match the source."
}
\end{verbatim}


\section{Hyperparameters}
\label{app:hyperparams}

\begin{table}[h]
\centering
\caption{Training hyperparameters for SFT and GRPO phases.}
\small
\begin{tabular}{lcc}
\toprule
\textbf{Hyperparameter} & \textbf{SFT} & \textbf{GRPO} \\
\midrule
Base model & \multicolumn{2}{c}{Qwen2.5-3B-Instruct} \\
Learning rate & $2 \times 10^{-5}$ & $2 \times 10^{-6}$ \\
LR scheduler & Cosine & Cosine \\
Warmup ratio & 0.05 & 0.05 \\
Weight decay & 0.01 & 0.01 \\
Epochs / Steps & 3 epochs & 350 steps \\
Batch size & 4 & --- \\
Gradient accumulation & 4 & --- \\
Max sequence length & 1280 & 512 (response) \\
Group size ($G$) & --- & 8 \\
Temperature & --- & 1.2 \\
Gradient clipping & --- & 0.1 \\
GPU memory utilization & --- & 0.4 \\
Precision & bf16 & bf16 \\
Hardware & \multicolumn{2}{c}{2$\times$ RTX 6000 Ada (48GB)} \\
\bottomrule
\end{tabular}
\end{table}

\end{document}
